\section{Related Work}
\label{sec:related-work}

\subsection{DisasterM3: the primary basis for this proposal}
\label{sec:related-disasterm3}

\emph{DisasterM3} \fcite{wang2025disasterm3} is, to date, the most comprehensive
attempt to benchmark VLMs on disaster-related RS vision tasks, and is the
direct basis for this proposal. It contributes 26,988 bi-temporal (pre/post)
satellite images and 123k instruction pairs spanning 36 historical disaster
events across 5 continents, grouped into 10 natural and man-made disaster
types (earthquakes, tsunamis, floods, explosions, storms, landslides,
tornadoes, conflicts, wildfires, volcanic eruptions). Imagery is drawn from
Maxar's WorldView series (optical) and from Capella Space and Umbra (SAR),
combined because extreme weather during disasters routinely prevents optical
sensors from imaging the scene at all. On top of this imagery, DisasterM3
defines nine tasks that progress from disaster-bearing-body and disaster-type
recognition, through damaged-building counting, damaged-road estimation, and
referring segmentation, to damaged-object relational reasoning and long-form
report/restoration-advice generation i.e. detection-, segmentation-, and
classification-style tasks alongside open-ended reasoning, unified under a
single instruction-following interface.

DisasterM3 evaluates 14 VLMs: five generic open models (LLaVA-OneVision,
InternVL3, Qwen2.5-VL, Kimi-VL, LLaVA-1.5), three commercial models (GPT-4o,
Claude-3, GPT-4.1), and seven RS-specific models (GeoChat, TeoChat, EarthDial,
GeoPixel, LISA, PSALM, HyperSeg); it then fine-tunes four of them
(Qwen2.5-VL-7B, InternVL3-8B, LISA-7B, PSALM-1.3B) on the dataset. The
headline findings are that (i) state-of-the-art VLMs struggle across the
board absent a disaster-specific corpus, (ii) a \emph{cross-sensor gap}
persists even after fine-tuning e.g. fine-tuned PSALM improves referring
segmentation by 40.8\% optical$\to$optical but only 23.7\% optical$\to$SAR --
and (iii) models show a density-dependent \emph{damage-counting} bias
(confident at the extremes of building density, unreliable in the mid-range),
which fine-tuning can shift toward an overfitting dilemma rather than resolve.
Crucially for this proposal, DisasterM3 reports the cross-sensor gap only as
a binary optical-vs-SAR split: because its SAR imagery is pooled from two
distinct commercial providers (Capella and Umbra) with different sensor
characteristics, and its optical imagery is exclusively Maxar, the paper
cannot and does not say whether the ``cross-sensor gap'' is really a
cross-\emph{provider} gap, a cross-\emph{resolution} gap, or genuinely a
modality (SAR vs.\ optical) effect.

\subsection{Prior work that DisasterM3 builds on}
\label{sec:related-prior}

DisasterM3 is not a first attempt so much as a merge of two threads pursued
by an overlapping set of authors. \emph{EarthVQA} \fcite{wang2024earthvqa},
led by the same first author (Junjue Wang) roughly a year and a half earlier,
introduced relational-reasoning visual question answering for RS imagery:
6,000 images of Wuhan, Changzhou, and Nanjing with 208,593 QA pairs spanning
judging, counting, and comprehensive relational-reasoning questions about
urban and rural governance (e.g.\ road/water renovation needs), plus the SOBA
(Semantic OBject Awareness) framework for answering them. EarthVQA is not
disaster-specific -- its objects of interest are governance-relevant urban
features, not damage -- but its task taxonomy (judging, counting, relational
reasoning, comprehensive analysis) is structurally the template that
DisasterM3's own task suite (damaged-object relational reasoning, damaged-
building counting, comprehensive report generation) generalizes from urban
planning to disaster response. Separately, \emph{BRIGHT}
\fcite{chen2025bright}, posted about four months before DisasterM3 and sharing
six of DisasterM3's authors,
built a building-damage-assessment dataset from exactly the three imagery
providers DisasterM3 later reuses: Maxar (optical) and Capella and Umbra
(SAR), across 14 regions at 0.3--1\,m resolution, evaluated with seven models
and adopted as the official dataset of the 2025 IEEE GRSS Data Fusion
Contest. DisasterM3's multi-sensor optical+SAR design is thus a direct
continuation of BRIGHT's data-collection pipeline, generalized from
single-task damage classification to a nine-task instruction-following VLM
benchmark but the provider-level distinctions BRIGHT's own collection
process preserves (three separately sourced providers) are exactly what get
collapsed into a single ``SAR'' pool once the data is repurposed inside
DisasterM3.

Two older, widely used datasets sit further back in this lineage and recur
throughout the wider survey (Section~\ref{sec:dataset-landscape}). \emph{xBD}
\cite{gupta2019xbd} established the pre/post, ordinal (no damage / minor /
major / destroyed) building-damage-classification paradigm on Maxar imagery
across 15 countries and remains the largest damage-labeled corpus available,
underpinning the classification-style task family that DisasterM3, BRIGHT,
and most datasets in our survey inherit. \emph{FloodNet}
\cite{rahnemoonfar2021floodnet} showed, four years before DisasterM3, that a
single detect/segment/classify/VQA task bundle could be built around one
disaster type (post-hurricane flooding) from UAV imagery validating the
``one model, several CV tasks, natural-language interface'' formulation at
centimeter scale well before VLM-scale, multi-hazard benchmarks existed.

\subsection{Concurrent work: MONITRS}
\label{sec:related-concurrent}

\emph{MONITRS} \fcite{revankar2025monitrs} (Cornell and Columbia; no author
overlap with the Wang/Yokoya lineage above) was posted on arXiv about eight
weeks after DisasterM3 and, like DisasterM3, was accepted to the NeurIPS 2025
Datasets and Benchmarks Track as a Spotlight. It targets the same
underlying problem from a different angle: rather than curating structured,
per-image instruction pairs, it pairs roughly 10,000 FEMA-catalogued disaster
events with temporal satellite-imagery sequences and natural-language
annotations mined from contemporaneous news coverage, and evaluates event-
level temporal understanding (e.g.\ correctly localizing when a storm's
impact began, which baseline MLLMs miss but MONITRS-fine-tuned models catch).
Two independent teams converging on ``benchmark and fine-tune VLMs on
disaster RS imagery'' within the same publication cycle is itself evidence
that this is a live, active problem rather than a solved one. Because
MONITRS is organized around event-level temporal/textual narrative rather
than DisasterM3's per-image structured task suite, the two datasets are
complementary axes for a broader evaluation rather than redundant ones. (An
earlier hypothesis behind this proposal was that \emph{EarthVQA} might
likewise have appeared in parallel with DisasterM3; the record instead shows
EarthVQA preceding it by roughly 18 months and sharing an author, which is
why it is treated above as prior lineage rather than as concurrent work.)
